\subsubsection{强最大值原理}\label{subsec:parabolic-strong-maximum}

本节的目标是以下结果.

\begin{Theorem}
\label{thm:parabolic-strong-maximum}
假设 $Lu\geq0$ ($Lu\leq0$). 令 $M=\sup_Du$ ($M=\inf_Du$). 假设在点 $(\bm{x}_0,t_0)\in D$ 处 $u=M$, 并且以下条件之一成立:
\begin{enumerate}
\item $c=0$, 而 $M$ 任意,
\item $c\leq0$ 且 $M\geq0$ ($M\leq0$),
\item $M=0$, 而 $c$ 任意.
\end{enumerate}
那么在 $\overline\Omega\times[0,t_0]$ 上 $u=M$.
\end{Theorem}

证明由三个引理得出. 从现在起, 始终假设 $Lu\geq0$ 而不再明确陈述; 显然, 以 $-u$ 代替 $u$ 即得 $Lu\leq0$ 的情形. 以下引理直接由 Lemma~\ref{lem:elliptic-boundary-point} 和 Remark~\ref{rem:boundary-point-semidefinite} 得出.

\begin{Lemma}
\label{lem:parabolic-ball-extreme-time}
设 $B\subset\mathbb{R}^{n+1}$ 为球, 满足 $\overline B\subset D$, 并假设 $B$ 内 $u<M$, 且 $u(\bm{x}_0,t_0)=M$, 其中 $(\bm{x}_0,t_0)\in\partial B$. 那么 $t_0$ 是 $t$ 在 $\overline B$ 中取得的最小值或最大值.
\end{Lemma}

接下来证明, 若在 $Q$ 内某点 $(\bm{x}_0,t_0)$ 处 $u<M$, 则在整个 $\Omega\times\{t_0\}$ 上 $u<M$.

\begin{Lemma}
\label{lem:parabolic-slice-strictness}
假设 $u(\bm{x}_0,t_0)<M$, 其中 $\bm{x}_0\in\Omega$ 且 $t_0\in(0,T)$. 那么对每个 $\bm{x}\in\Omega$, 都有 $u(\bm{x},t_0)<M$.
\end{Lemma}

\begin{Proof}
假设 $\Omega\times\{t_0\}$ 包含使 $u=M$ 的点. 那么可以在 $\Omega$ 中找到点 $\bm{x}_1$ 和 $\bm{x}_2$, 使 $u(\bm{x}_1,t_0)=M$, $u(\bm{x}_2,t_0)<M$, 且连接 $(\bm{x}_1,t_0)$ 与 $(\bm{x}_2,t_0)$ 的线段 $L$ 完全位于 $\Omega\times\{t_0\}$ 内. 移动点 $\bm{x}_1$, 总可以使沿 $L$ 有 $u<M$. 现在令
\begin{equation}\label{eq:parabolic-slice-distance-radius}
\delta=\min(|\bm{x}_1-\bm{x}_2|,\operatorname{dist}(L,\partial D)).
\end{equation}
对满足 $0<|\bm{x}-\bm{x}_1|<\delta$ 的 $\bm{x}\in L$, 令
\begin{equation}\label{eq:parabolic-slice-maximum-set-distance}
d(\bm{x})=\operatorname{dist}((\bm{x},t_0),Q\cap\{u=M\}).
\end{equation}
显然, $d(\bm{x})\leq|\bm{x}-\bm{x}_1|$.

由 Lemma~\ref{lem:parabolic-ball-extreme-time}, $Q\cap\{u=M\}$ 中离 $(\bm{x},t_0)$ 最近的点具有 $(\bm{x},t)$ 的形式, 因而有 $u(\bm{x},t_0+d(\bm{x}))=M$ 或 $u(\bm{x},t_0-d(\bm{x}))=M$. 设 $\bm{n}$ 为沿 $L$ 方向的单位向量. 当 $|\epsilon|>0$ 充分小时, $d(\bm{x}+\epsilon\bm{n})$ 有定义, 并由 Pythagorean 定理有
\begin{equation}\label{eq:parabolic-slice-distance-forward-bound}
d(\bm{x}+\epsilon\bm{n})\leq\sqrt{\epsilon^2+d(\bm{x})^2}.
\end{equation}
由相同论证,
\begin{equation}\label{eq:parabolic-slice-distance-reverse-bound}
d(\bm{x})\leq\sqrt{\epsilon^2+d(\bm{x}+\epsilon\bm{n})^2},
\end{equation}
因而
\begin{equation}\label{eq:parabolic-slice-distance-lower-bound}
d(\bm{x}+\epsilon\bm{n})\geq\sqrt{-\epsilon^2+d(\bm{x})^2}.
\end{equation}
由 \eqref{eq:parabolic-slice-distance-reverse-bound} 和 \eqref{eq:parabolic-slice-distance-lower-bound}, 得到 $d(\bm{x}+\epsilon\bm{n})$ 在 $\epsilon=0$ 处的导数存在且为零. 因此 $d(\bm{x})$ 沿 $L\cap\{0<|\bm{x}-\bm{x}_1|<\delta\}$ 为常数. 但这产生矛盾, 因为 $d(\bm{x})\ne0$, 而当 $\bm{x}\to\bm{x}_1$ 时 $d(\bm{x})\to0$.
\end{Proof}

所需的最后一个引理如下.

\begin{Lemma}
\label{lem:parabolic-time-strictness}
设 $0\leq t_0<t_1\leq T$, 并假设在 $\Omega\times(t_0,t_1)$ 内 $u<M$. 那么在 $\Omega\times\{t_1\}$ 上 $u<M$.
\end{Lemma}

\begin{Proof}
反设存在 $\bm{x}_1\in\Omega$, 满足 $u(\bm{x}_1,t_1)=M$. 定义
\begin{equation}\label{eq:parabolic-time-propagation-barrier}
v(\bm{x},t)=\exp\bigl(-|\bm{x}-\bm{x}_1|^2-\alpha(t-t_1)\bigr)-1,
\end{equation}
其中 $\alpha>0$ 取得很大. 计算得到\[
\begin{aligned}
Lv(\bm{x},t)&=\exp\bigl(-|\bm{x}-\bm{x}_1|^2-\alpha(t-t_1)\bigr)\\
&\quad\times\bigl(4a_{ij}(x_i-x_{1i})(x_j-x_{1j})-2(a_{ii}+b_i(x_i-x_{1i}))+\alpha\bigr)+cv.
\end{aligned}
\]若 $N$ 为 $(\bm{x}_1,t_1)$ 的一个小邻域, 则只要将 $\alpha$ 取得足够大, 就有在 $N\cap\{t\leq t_1\}$ 内 $Lv>0$. 现在设 $A$ 为抛物面 $|\bm{x}-\bm{x}_1|^2+\alpha(t-t_1)=0$ 和以 $(\bm{x}_1,t_1)$ 为球心的充分小的球面所围成的区域. $A$ 的边界分为两部分, 一部分在抛物面上, 那里 $v=0$, 另一部分在球面上且在抛物面内部, 那里 $u<M$.

因此, 若将 $\epsilon>0$ 取得很小, 就有在 $\partial A$ 上 $u+\epsilon v-M\leq0$, 且在 $A$ 内 $L(u+\epsilon v-M)\geq-cM\geq0$.

若 $c\leq0$, 则弱最大值原理表明, 在整个 $A$ 内 $u+\epsilon v-M\leq0$. 在点 $(\bm{x}_1,t_1)$ 处取关于 $t$ 的导数, 得到
\begin{equation}\label{eq:parabolic-terminal-time-derivative-bound}
u_t+\epsilon v_t=u_t-\epsilon\alpha\geq0,
\end{equation}
即 $u_t>0$. 但由于 $(\bm{x}_1,t_1)$ 是关于 $\bm{x}$ 的最大值点, 得到在该点
\begin{equation}\label{eq:parabolic-terminal-operator-bound}
Lu\leq-u_t+cu\leq-u_t,
\end{equation}
又因假设了 $Lu\geq0$, 得出 $u_t\leq0$, 矛盾.

若 $M=0$, 则可如 Lemma~\ref{lem:elliptic-boundary-point} 的证明中那样, 以 $(L-c^+)$ 代替 $L$, 应用相同论证.
\end{Proof}

现在定理的证明已清楚. 对每个 $t>0$, 或者对所有 $\bm{x}\in\Omega$ 都有 $u(\bm{x},t)=M$, 或者对所有这些点都有 $u(\bm{x},t)<M$ (Lemma~\ref{lem:parabolic-slice-strictness}). 所有使 $u<M$ 的 $t$ 构成的集合 $Z$ 是开的, 因而是可数个区间的并. 由 Lemma~\ref{lem:parabolic-time-strictness}, 这些区间都不能有上端点. 唯一的可能性就是对某个 $t_0$, 有 $Z=(t_0,T]$, 这正是定理的结论.

最后指出, 若 $\partial\Omega$ 属于 $C^2$ 类, 且 $u$ 直到边界都为 $C^1$ 的, 则 Lemma~\ref{lem:elliptic-boundary-point} 适用, 即若 $\bm{x}_0\in\partial\Omega$, $t\in(0,T)$, $u(\bm{x}_0,t)=M$, 但对 $\bm{x}\in\Omega$ 有 $u(\bm{x},t)<M$, 则 $\partial u/\partial n(\bm{x}_0,t)>0$.

\paragraph*{Problems}

\begin{enumerate}[label=\thesection.\arabic*.,ref=\thesection.\arabic*,start=22,series=book-problems]

\item\label{ex:parabolic-neumann-comparison} 假设 $\Omega$ 有界, $\partial\Omega$ 属于 $C^2$ 类, 且 $u,v\in C^2(D)\cap C^1(\overline D)$. 再假设对 $(\bm{x},t)\in D$ 有 $Lu\leq Lv$, 对 $\bm{x}\in\Omega$ 有 $u(\bm{x},0)\geq v(\bm{x},0)$, 并且对 $(\bm{x},t)\in\partial\Omega\times(0,T)$ 有 $\partial u/\partial n\geq\partial v/\partial n$. 证明在 $D$ 内 $u\geq v$.

\item\label{ex:parabolic-stable-equilibrium-nonnegative} 假设 $\Omega$ 有界, $f\leq0$, 且 $f$ 及 $L$ 的系数均不依赖于 $t$. 考虑方程 $Lu=f$, 边界条件为在 $\partial\Omega\times(0,\infty)$ 上 $u=0$. 平衡解是一个不依赖于时间的解. 若对每个满足在 $\partial\Omega$ 上 $u_0=0$ 的连续函数 $u_0(\bm{x})$, 都存在解 $u$, 满足在 $\Omega\times(0,\infty)$ 内 $Lu=f$, 在 $\partial\Omega\times(0,\infty)$ 上 $u=0$, 对 $\bm{x}\in\Omega$ 有 $u(\bm{x},0)=u_0(\bm{x})$, 并且当 $t\to\infty$ 时 $u(\bm{x},t)\to v(\bm{x})$, 则称平衡解 $v$ 稳定. 证明稳定的平衡解必为非负.

\item\label{ex:whole-space-bounded-heat-maximum} 令 $D=\mathbb{R}^n\times(0,T]$. 设 $u\in C^2(D)\cap C(\overline D)$ 为热方程 $u_t=\Delta u$ 的有界解. 证明 $\sup_Du\leq\sup_{\mathbb{R}^n}u(\cdot,0)$. Hint: 考虑函数 $v=u-\epsilon(2nt+|\bm{x}|^2)$.

\item\label{ex:semilinear-parabolic-uniqueness} 设 $\Omega$ 有界, 且 $f:\mathbb{R}\to\mathbb{R}$ 属于 $C^1$ 类. 设 $u_0\in C(\overline\Omega)$, 在 $\partial\Omega$ 上满足 $u_0=0$. 证明对任意 $T>0$, 方程 $u_t=\Delta u+f(u)$ 在边界条件为 $\partial\Omega$ 上 $u=0$, 初始条件为 $u(\bm{x},0)=u_0(\bm{x})$ 时, 在 $\Omega\times(0,T)$ 上至多有一个解.

\item\label{ex:reaction-diffusion-positivity} 设 $\Omega$ 有界, 且 $f,g:\mathbb{R}^2\to\mathbb{R}$ 属于 $C^1$ 类. 设 $u$, $v$ 在 $\Omega\times(0,T)$ 上满足方程 $u_t=\Delta u+f(u,v)$, $v_t=\Delta v+g(u,v)$, 边界条件为 $\partial\Omega$ 上 $u=v=0$. 假设对每个 $v$ 都有 $f(0,v)\geq0$, 且对每个 $u$ 都有 $g(u,0)\geq0$. 证明, 若对每个 $\bm{x}\in\Omega$ 都有 $u(\bm{x},0)\geq0$, $v(\bm{x},0)\geq0$, 则 $u\geq0$, $v\geq0$. 这里描述的这类系统出现在反应扩散问题和数学生物学中, 其中 $u$ 与 $v$ 表示物种的浓度或种群数量.

\end{enumerate}
